
Our work intersects three research areas: sampling-based uncertainty quantification for LLMs, hallucination detection benchmarks, and the behavioral effects of RLHF fine-tuning on LLM output distributions.

\subsubsection{Sampling-Based Consistency for Uncertainty Quantification}

The use of sampling-based consistency as an uncertainty signal in language models was established by Wang et al. [2022] in the context of chain-of-thought reasoning: majority vote consistency over multiple samples implicitly captures model confidence. Kuhn et al. [2023] formalized this intuition as \textit{Semantic Uncertainty} --- using NLI to cluster semantically equivalent answers and computing entropy over these semantic equivalence classes, rather than over token sequences. Applied to TriviaQA and NaturalQuestions with GPT-3 and smaller models on open-ended generation, Semantic Uncertainty achieves meaningful AUROC and outperforms token-probability baselines. Critically, the models evaluated by Kuhn et al. are base or lightly fine-tuned models where the stochastic hallucination assumption naturally holds.

Manakul et al. [2023] introduced SelfCheckGPT, which directly operationalizes NLI-based pairwise consistency as a hallucination detector. Evaluated on WikiBio biography generation with GPT-3, SelfCheckGPT achieves AUROC ~0.65--0.75, confirming that consistent outputs signal factual accuracy in open-ended generation. SelfCheckGPT is the direct methodological predecessor of our SMC-NLI scorer.

Lin et al. [2024] conducted a systematic comparison of black-box uncertainty methods, finding that sampling-based consistency is competitive with white-box (logit-based) methods when logits are unavailable. Xiong et al. [2023] showed that sampling-based UQ generalizes better across model families than verbalized confidence.

\textbf{The gap:} All prior work establishing sampling-based consistency validates on base or lightly-tuned models on open-ended generation tasks. None has systematically evaluated whether this paradigm holds for RLHF-fine-tuned models on structured factual QA. We fill this gap.

\subsubsection{Hallucination Detection Benchmarks}

Hallucination detection has been studied across multiple benchmark types. TruthfulQA \citep{Lin2022TruthfulQA} evaluates models' propensity to repeat imitative falsehoods. TriviaQA \citep{Joshi2017} and NaturalQuestions \citep{Kwiatkowski2019} provide ground-truth factual QA with exact-match evaluation, enabling model-specific hallucination assessment.

HaluEval \citep{Li2023HaluEval}, which we use as our primary benchmark, was constructed by sampling ChatGPT responses and asking ChatGPT to generate hallucinated versions, producing binary correct/hallucinated labels for QA, dialogue, and summarization tasks. HaluEval has been widely adopted due to its scale and accessibility.

However, HaluEval's construction protocol introduces a cross-model evaluation concern: the binary labels reflect ChatGPT's hallucination patterns, not those of the model being evaluated. A model that answers correctly where ChatGPT hallucinated --- or vice versa --- will receive inaccurate labels, making AUROC evaluation ambiguous. Our results suggest this is a non-trivial concern for Llama-3-8B-Instruct. Benchmarks with model-agnostic ground truth (TriviaQA exact match, NQ F1) avoid this confound and should be preferred for cross-model hallucination detection evaluation.

\subsubsection{Effects of RLHF Fine-Tuning on Model Output Distributions}

RLHF training produces models that are more deterministic and consistent in their outputs [Ziegler et al., 2019; Ouyang et al., 2022] --- a desirable property for deployment but potentially problematic for uncertainty quantification. Instruction-following fine-tuning tends to sharpen model output distributions \citep{Bai2022Training}. At temperature=0.7, an instruction-tuned model may produce a much narrower distribution over output tokens than an equivalently-sized base model, because RLHF training has reinforced specific response patterns. If this sharpening is strong enough, even ''uncertain'' responses may be generated consistently across samples, eliminating the consistency-based discriminative signal.

This theoretical concern has not been empirically tested against sampling-based hallucination detectors. To our knowledge, our work provides the first direct empirical evidence of this effect: Llama-3-8B-Instruct produces nearly identical consistency scores for correctly-labeled (mean SMC-NLI=0.6236) and hallucinated-labeled (mean SMC-NLI=0.6299) questions, with a gap of 0.006.

\subsubsection{Positioning Our Work}

Our contribution is complementary to, not contradictory of, prior literature. SelfCheckGPT and Semantic Uncertainty establish that sampling-based consistency works in the stochastic hallucination regime (open-ended generation, base/lightly-tuned models). We establish that instruction-tuned models on structured factual QA exhibit a different regime (systematic confabulation) where these methods fail. The two findings together define a regime boundary that practitioners need to know before deploying consistency-based detectors.

